\documentclass[12pt, letterpaper]{article}

\usepackage[utf8]{inputenc}
\usepackage[T1]{fontenc}
\usepackage[spanish]{babel}
\usepackage{geometry}
\usepackage{amsmath, amssymb, amsthm}
\usepackage{booktabs}
\usepackage{array}
\usepackage{microtype}
\usepackage{setspace}
\usepackage{parskip}
\usepackage{titlesec}
\usepackage{fancyhdr}
\usepackage{enumitem}
\usepackage{bm}
\usepackage{mathtools}
\usepackage{tcolorbox}
\usepackage{hyperref}

\tcbuselibrary{skins, breakable}

\geometry{top=2.8cm, bottom=3.0cm, left=3.2cm, right=3.2cm, headheight=14pt}
\setstretch{1.20}

\pagestyle{fancy}
\fancyhf{}
\fancyhead[C]{\small\textsc{Econometría --- Gujarati --- Ejercicio 13.12}}
\fancyfoot[C]{\small\thepage}
\renewcommand{\headrulewidth}{0.4pt}
\renewcommand{\footrulewidth}{0pt}

\titleformat{\section}
  {\large\bfseries\scshape}{\thesection.}{0.6em}{}
  [\vspace{-4pt}\rule{\linewidth}{0.4pt}]
\titleformat{\subsection}{\normalsize\bfseries}{\thesubsection.}{0.5em}{}
\titleformat{\subsubsection}{\normalsize\itshape}{}{0em}{}

\newtcolorbox{clave}{
  colback=white, colframe=black,
  boxrule=0.5pt, arc=3pt,
  left=8pt, right=8pt, top=6pt, bottom=6pt,
  breakable
}

\newcommand{\E}{\operatorname{E}}
\newcommand{\Cov}{\operatorname{Cov}}
\newcommand{\Var}{\operatorname{Var}}

\begin{document}

\begin{center}
  {\LARGE\bfseries Sesgo por Variable Omitida en el Intercepto}\\[6pt]
  {\LARGE\bfseries Demostración de $\E(\hat{\alpha}_1) = \beta_1 + \beta_3(\bar{X}_3 - b_{32}\bar{X}_2)$}\\[12pt]
  {\large\itshape Especificación incorrecta por omisión de una variable relevante:}\\
  {\large\itshape efecto sobre el estimador MCO del intercepto}\\[14pt]
  \rule{0.55\linewidth}{0.6pt}\\[8pt]
  {\normalsize Ejercicio~13.12 --- \textit{Econometría}, Gujarati \& Porter, 5.a ed.}\\[4pt]
  {\small\textsc{Demostración conceptual, teórica y algebraica}}
\end{center}

\vspace{10pt}

%------------------------------------------------------------------------------
\section{Marco teórico y planteamiento}
%------------------------------------------------------------------------------

El ejercicio estudia las consecuencias de \textbf{omitir una variable relevante}
sobre las propiedades del estimador MCO del intercepto. Si bien el sesgo
del coeficiente de pendiente es el resultado más conocido en la literatura
sobre errores de especificación, el intercepto también resulta afectado de
manera sistemática.

\subsection{Modelo verdadero (PRF)}

Supóngase que el proceso generador de datos verdadero es el modelo de
regresión múltiple:
\begin{equation}
  Y_i = \beta_1 + \beta_2 X_{2i} + \beta_3 X_{3i} + u_i,
  \qquad i = 1, \ldots, n,
  \label{eq:verdadero}
\end{equation}
donde $u_i$ satisface las propiedades clásicas: $\E[u_i] = 0$,
$\Var(u_i) = \sigma^2$, y los regresores son fijos (no estocásticos) en
muestreos repetidos. Los parámetros $\beta_2$ y $\beta_3$ son
\emph{ambos distintos de cero}, por lo que $X_{3i}$ es una
\textbf{variable relevante}.

\subsection{Modelo estimado (SRF) con variable omitida}

En la práctica, el investigador omite $X_3$ y estima en su lugar el
modelo restringido:
\begin{equation}
  Y_i = \alpha_1 + \alpha_2 X_{2i} + v_i,
  \label{eq:omitido}
\end{equation}
donde el error compuesto $v_i$ recoge tanto la perturbación original como
la influencia de la variable excluida:
\[
  v_i = \beta_3 X_{3i} + u_i.
\]
Conceptualmente, $\hat{\alpha}_1$ y $\hat{\alpha}_2$ ya no son estimadores
del verdadero $\beta_1$ y $\beta_2$, pues el supuesto de exogeneidad
$\E[v_i \mid X_{2i}] = 0$ se viola siempre que $X_{3i}$ esté correlacionada
con $X_{2i}$.

%------------------------------------------------------------------------------
\section{Resultado auxiliar: sesgo del coeficiente de pendiente (ec.\ 13.3.3)}
%------------------------------------------------------------------------------

Antes de demostrar el resultado pedido, se recuerda el sesgo del estimador
de la pendiente, derivado en la ecuación (13.3.3) del texto:
\begin{equation}
  \E(\hat{\alpha}_2) = \beta_2 + \beta_3\, b_{32},
  \label{eq:sesgo_pendiente}
\end{equation}
donde $b_{32}$ es el coeficiente de pendiente de la \textbf{regresión auxiliar}
de la variable omitida $X_3$ sobre la variable incluida $X_2$:
\begin{equation}
  X_{3i} = a + b_{32}\, X_{2i} + \text{residuo}.
  \label{eq:reg_auxiliar}
\end{equation}
El término $b_{32}$ captura la relación lineal entre $X_3$ y $X_2$;
si las dos variables son ortogonales, $b_{32} = 0$ y el sesgo en
$\hat{\alpha}_2$ desaparece. Este resultado auxiliar se utilizará en el
Paso~4 de la demostración principal.

%------------------------------------------------------------------------------
\section{Demostración: $\E(\hat{\alpha}_1) = \beta_1 + \beta_3(\bar{X}_3 - b_{32}\bar{X}_2)$}
%------------------------------------------------------------------------------

La demostración se desarrolla en cinco pasos secuenciales.

\subsection{Paso 1 --- Expresión del estimador MCO del intercepto}

En cualquier regresión simple de la forma~\eqref{eq:omitido}, el estimador
MCO del intercepto se obtiene de la condición de primer orden que impone
que la recta de regresión muestral pase por el punto de medias
$(\bar{X}_2,\, \bar{Y})$:
\begin{equation}
  \hat{\alpha}_1 = \bar{Y} - \hat{\alpha}_2\,\bar{X}_2.
  \label{eq:intercepto_mco}
\end{equation}
Esta identidad es exacta en cualquier muestra finita y no depende de
ningún supuesto distribucional.

\subsection{Paso 2 --- Media muestral de $Y$ a partir del modelo verdadero}

Promediando ambos lados de la ecuación~\eqref{eq:verdadero} sobre las
$n$ observaciones:
\begin{equation}
  \bar{Y} = \beta_1 + \beta_2\,\bar{X}_2 + \beta_3\,\bar{X}_3 + \bar{u},
  \label{eq:media_Y}
\end{equation}
donde $\bar{u} = n^{-1}\sum_{i=1}^n u_i$ es la media muestral de las
perturbaciones.

\subsection{Paso 3 --- Sustitución y aplicación del operador de esperanza}

Sustituyendo~\eqref{eq:media_Y} en~\eqref{eq:intercepto_mco}:
\[
  \hat{\alpha}_1
  = \bigl(\beta_1 + \beta_2\,\bar{X}_2 + \beta_3\,\bar{X}_3 + \bar{u}\bigr)
    - \hat{\alpha}_2\,\bar{X}_2.
\]
Aplicando el operador de valor esperado a ambos lados, y usando que los
regresores son fijos (por lo que $\bar{X}_2$ y $\bar{X}_3$ son
constantes), $\beta_1$, $\beta_2$, $\beta_3$ son parámetros fijos,
y $\E[\bar{u}] = 0$:
\begin{align}
  \E(\hat{\alpha}_1)
  &= \beta_1 + \beta_2\,\bar{X}_2 + \beta_3\,\bar{X}_3
     + \underbrace{\E[\bar{u}]}_{=\,0}
     - \E(\hat{\alpha}_2)\,\bar{X}_2 \notag\\[4pt]
  &= \beta_1 + \beta_2\,\bar{X}_2 + \beta_3\,\bar{X}_3
     - \E(\hat{\alpha}_2)\,\bar{X}_2.
  \label{eq:paso3}
\end{align}

\subsection{Paso 4 --- Sustitución del sesgo conocido de la pendiente}

Reemplazamos $\E(\hat{\alpha}_2)$ en~\eqref{eq:paso3} usando el
resultado auxiliar~\eqref{eq:sesgo_pendiente}:
\begin{align}
  \E(\hat{\alpha}_1)
  &= \beta_1 + \beta_2\,\bar{X}_2 + \beta_3\,\bar{X}_3
     - (\beta_2 + \beta_3\,b_{32})\,\bar{X}_2.
  \label{eq:paso4}
\end{align}

\subsection{Paso 5 --- Simplificación algebraica}

Distribuyendo $\bar{X}_2$ en el último término de~\eqref{eq:paso4}:
\begin{align}
  \E(\hat{\alpha}_1)
  &= \beta_1
     + \beta_2\,\bar{X}_2
     + \beta_3\,\bar{X}_3
     - \beta_2\,\bar{X}_2
     - \beta_3\,b_{32}\,\bar{X}_2.
  \label{eq:paso5a}
\end{align}
Los términos $\beta_2\,\bar{X}_2$ y $-\beta_2\,\bar{X}_2$ se cancelan:
\begin{align}
  \E(\hat{\alpha}_1)
  &= \beta_1 + \beta_3\,\bar{X}_3 - \beta_3\,b_{32}\,\bar{X}_2.
  \label{eq:paso5b}
\end{align}
Finalmente, factorizando $\beta_3$:
\begin{equation}
  \boxed{
    \E(\hat{\alpha}_1) = \beta_1 + \beta_3\,(\bar{X}_3 - b_{32}\,\bar{X}_2)
  }
  \label{eq:resultado}
\end{equation}
que es exactamente lo que se quería demostrar. \hfill $\blacksquare$

%------------------------------------------------------------------------------
\section{Interpretación del resultado}
%------------------------------------------------------------------------------

\subsection{Naturaleza del sesgo}

La ecuación~\eqref{eq:resultado} muestra que el estimador $\hat{\alpha}_1$
es en general un estimador \textbf{sesgado e inconsistente} de $\beta_1$.
El sesgo es:
\begin{equation}
  \text{Sesgo}(\hat{\alpha}_1)
  = \E(\hat{\alpha}_1) - \beta_1
  = \beta_3\,(\bar{X}_3 - b_{32}\,\bar{X}_2).
  \label{eq:sesgo}
\end{equation}
Este sesgo tiene una interpretación geométrica precisa: el intercepto
estimado absorbe el \emph{efecto promedio} de la variable omitida
($\beta_3\,\bar{X}_3$), pero descontando la parte de ese efecto que ya
fue atribuida, de manera espuria, al coeficiente de pendiente a través
de la correlación entre $X_3$ y $X_2$ (capturada por $b_{32}\,\bar{X}_2$).

\subsection{Condiciones de insesgadez}

El intercepto $\hat{\alpha}_1$ sería un estimador insesgado de $\beta_1$
únicamente si el sesgo~\eqref{eq:sesgo} se anulara. Esto ocurre en
alguno de los siguientes casos:

\begin{enumerate}[leftmargin=1.8cm, label=\textbf{\arabic*.}]
  \item $\beta_3 = 0$: la variable $X_3$ no pertenece realmente al
    modelo, por lo que omitirla no genera sesgo alguno.
  \item $\bar{X}_3 = b_{32}\,\bar{X}_2$: la media muestral de la
    variable omitida coincide exactamente con la predicción lineal de
    $X_3$ a partir de $X_2$, evaluada en la media de $X_2$. Esta
    condición es algebraicamente posible pero estadísticamente
    improbable en datos económicos reales.
  \item $\bar{X}_3 = 0$ y $b_{32} = 0$ simultáneamente: ambas
    variables tienen media cero y son ortogonales, situación que puede
    lograrse artificialmente centrando las variables pero que no
    elimina el sesgo sobre la pendiente.
\end{enumerate}

\subsection{Relación con el sesgo de la pendiente}

Comparando~\eqref{eq:resultado} con~\eqref{eq:sesgo_pendiente}, se aprecia
la asimetría del problema: el sesgo de $\hat{\alpha}_2$ depende de
$b_{32}$ (la \emph{relación} entre las variables), mientras que el sesgo
de $\hat{\alpha}_1$ depende de $(\bar{X}_3 - b_{32}\,\bar{X}_2)$ (una
combinación de \emph{niveles} y \emph{relación}). Así, incluso si
$b_{32} = 0$ (variables ortogonales, lo que elimina el sesgo de la
pendiente), el intercepto seguirá sesgado siempre que $\bar{X}_3 \neq 0$.

\begin{clave}
  \textbf{Resumen:} La omisión de una variable relevante $X_3$ produce:
  \begin{itemize}[leftmargin=1.2cm]
    \item $\E(\hat{\alpha}_2) = \beta_2 + \beta_3\,b_{32}$ \quad
      (sesgo en la pendiente, ec.\ 13.3.3).
    \item $\E(\hat{\alpha}_1) = \beta_1 + \beta_3\,(\bar{X}_3 - b_{32}\,\bar{X}_2)$
      \quad (sesgo en el intercepto, resultado demostrado).
  \end{itemize}
  Ambos estimadores son \textbf{sesgados e inconsistentes}. El
  intercepto absorbe el efecto medio de la variable omitida corregido por
  la correlación con el regresor incluido. En general, no existe
  ningún remedio dentro del MCO: la única solución correcta es
  \textbf{incluir la variable omitida} en la especificación.
\end{clave}

\end{document}
